\chapter{Практика}
\label{ch:chap2}

% \begin{figure}[H]
%     \centering
%     \includegraphics{carrier.png}
%     \caption{Несущие колебания}
% \end{figure}


\section{Сигнал с шумом и его автокорреляция}

Пусть есть сигнал, который при прохождении через канал становится зашумленным. Его можно
записать как $X(t) = S(t) + N(t)$. \\

АКФ для такого сигнала можно записать как:

$$R_x(\tau) = E[S(t)S(t+\tau)] = E[(S(t)+N(t))(S(t+\tau)(N(t+\tau)))]$$

Раскроем скобки:

$$R_x(\tau) = E[S(t)S(t+\tau)] + E[S(t)N(t+\tau)] + E[N(t)S(t+\tau)] + E[N(t)N(t+\tau)]$$

Если предполагать, что шум и сигнал независимы, то $E[S(t)N(t+\tau)]$ и $E[N(t)S(t+\tau)]$ равны нулю,
получается, что:

$$R_x(\tau) = E[S(t)S(t+\tau)] + E[N(t)N(t+\tau)] = R_s(\tau) + R_n(\tau)$$

АКФ для шума $R_n(\tau)$ будет выглядеть (в идеале), как дельта функция. Будет один
пик при лаге равном 0. Таким образом, остается только $R_s(\tau)$. \\

Если сигнал периодичен и имеет какую-то структуру, то $R_s(\tau)$ будет похожа на сам сигнал.

\section{Задания}

Сгенерируем СП с нормальным и равномерным распределением. Также сгенерируем сигнал в виде
синусоиды и наложим на сигнал шум:

\begin{figure}[H]
    \centering
    \includegraphics{10.png}
    \caption{Сигнал и разные шумы}
\end{figure}

\begin{figure}[H]
    \centering
    \includegraphics{20.png}
    \caption{Зашумленный сигнал}
\end{figure}

Теперь вычислим АКФ:

\begin{figure}[H]
    \centering
    \includegraphics{30.png}
    \caption{АКФ сигнала с шумом}
\end{figure}

Можно увидеть, что шум исчез, а благодаря периодичности сигнала
АКФ похожа на исходный сигнал.\\

Повысим дисперсию нормального шума. Теперь шум по амплитуде больше сигнала.
В таком случае АКФ получится не такой гладкой, но она всё равно напоминает исходный сигнал


\begin{figure}[H]
    \centering
    \includegraphics{40.png}
    \caption{АКФ сигнала с сильным шумом}
\end{figure}

Случайный процесс состоит из суммы быстрых и медлленных изменений. Эти компоненты можно выделить с помощью ФНЧ 
или ФВЧ

\begin{figure}[H]
    \centering
    \includegraphics{50.png}
    \caption{ИХ фильтра}
\end{figure}

\begin{figure}[H]
    \centering
    \includegraphics{60.png}
    \caption{Быстрая и медленная компонента}
\end{figure}

Теперь посчитаем АКФ для этих процессов

\begin{figure}[H]
    \centering
    \includegraphics{70.png}
    \caption{АКФ для медленной и быстрой компоненты}
\end{figure}

\begin{figure}[H]
    \centering
    \includegraphics{80.png}
    \caption{АКФ для медленной и быстрой компоненты}
\end{figure}

Если приблизить АКФ, то можно заметить, что интервал корреляции
быстрой компоненты меньше, чем интервал корреляции медленной компоненты.
Это говорит о том, что по интервалу корреляции мы можем судить о поведении
процесса, а это очень важно, потому что по числовым характеристикам
процесса мы не можем судить о его поведении, а АКФ как раз это может показать.







\endinput 